% Conclusion

Biological Petri net theory lacks formal semantics for hierarchical regulatory signals consumed during decision-making, preventing expression of regulatory consumption phenomena essential for commitment finality. We introduced Signal Hierarchy Theory providing signal places ($\Psi$) with consumption semantics, enabling hierarchical preemption where higher regulatory layers control lower metabolic layers. The formalism bridges the expressiveness gap between metabolic mass transfer (normal arcs) and regulatory information flow (signal flow arcs) within unified Bio-PN framework.

\vspace{3pt}
\noindent The formalism enables analytical capabilities previously difficult to express in classical Bio-PN: threshold determination (energy concentrations triggering pathway commitment), basin geometry quantification (cell fate reliability via reachability analysis), and decision finality analysis (irreversible commitment through signal consumption). Proof-of-concept application to \textit{B. subtilis} sporulation demonstrates feasibility: the formalism computes ATP commitment threshold at 2.38 mM, matching experimental measurement (2.21 mM, 7\% error). This quantitative prediction illustrates advantages over classical Bio-PN lacking explicit consumption semantics for regulatory signals.

\vspace{3pt}
\noindent Researchers can now formally pose questions requiring signal consumption semantics: What energy threshold triggers irreversible pathway commitment? How reliable is cell fate convergence within attractor basins? When does regulatory decision become irreversible? Signal hierarchy provides formal semantics and mathematical methods for addressing these questions. Broader validation across diverse biological systems would strengthen evidence for the formalism's general applicability.

\vspace{3pt}
\noindent Potential applications span synthetic biology (engineered decision circuits with formal commitment properties), drug discovery (ATP-dependent regulatory hierarchies with computed threshold modulation), and precision medicine (quantitative cell fate prediction under therapeutic intervention). The unified framework bridges molecular mechanisms with cellular decision-making, potentially supporting rational design of biological systems with formal guarantees and predictive accuracy.

\vspace{3pt}
\noindent Signal Hierarchy Theory extends Bio-PN formalism with consumption semantics and hierarchical organization, capturing both metabolic detail and regulatory control in coherent theoretical framework. The \textit{B. subtilis} application demonstrates proof-of-concept for quantitative systems biology analyses. Open-source implementation ensures reproducibility and enables community evaluation of the approach.
